
\chapter{Problem Formulation}
\label{ch:problem-formulation}

This chapter formalises the system model, telemetry fault model, battery
dynamics, and safety constraints that underpin the ORIUS framework.  We
state the research questions mathematically and delineate the scope
boundary of the dissertation.

\section{ORIUS System Context}
\label{sec:orius-context}

ORIUS is a renewable-aware battery dispatch platform that operates at
the \emph{optimise-to-dispatch boundary}: an upstream forecast-and-optimise
pipeline produces a candidate dispatch action~$a_t^{\star}$, and a
downstream safety layer either certifies or corrects the action before it
reaches the physical plant.  The system architecture comprises:

\begin{enumerate}
  \item \textbf{Forecast stage.}  CQR models predict load, solar, and
    wind trajectories over a configurable horizon, producing both point
    forecasts and conformal prediction intervals.
  \item \textbf{Optimise stage.}  A deterministic LP (or a pluggable
    alternative) maximises arbitrage revenue subject to SOC, power, and
    ramp constraints evaluated on the \emph{observed} state.
  \item \textbf{Dispatch stage (DC3S).}  The safety layer scores
    observation quality, inflates uncertainty, tightens action bounds, and
    issues a safety certificate for the final dispatched action.
\end{enumerate}

The boundary between stages~2 and~3 is the \emph{optimise-to-dispatch
boundary}.  Everything upstream of this boundary operates on the observed
state~$\hat{x}_t$; the safety layer's task is to ensure that the
dispatched action is safe on the \emph{true} state~$x_t$ despite the gap
between the two.

\section{ORIUS-MDP Formalisation}
\label{sec:orius-mdp}

\begin{definition}[ORIUS-MDP]
\label{def:orius-mdp}
An ORIUS-controlled cyber-physical system is modelled as the 11-tuple
\[
  (\mathcal{S}, \mathcal{O}, \mathcal{W}, \mathcal{A}, \mathcal{C},
   \mathcal{T}, \mathcal{H}, \Omega, \mathcal{R}, \gamma, \mu_0),
\]
where $\mathcal{S}$ is the true-state space, $\mathcal{O}$ the observation
space, $\mathcal{W} = [w_{\min}, 1]$ the observation-quality space,
$\mathcal{A}$ the action space, and $\mathcal{C}\subsetneq\mathcal{S}$ the
closed safety set.  The true transition kernel is
$\mathcal{T}: \mathcal{S}\times\mathcal{A}\to\Delta(\mathcal{S})$, the
observation kernel is $\mathcal{H}: \mathcal{S}\to\Delta(\mathcal{O})$, the
quality map is $\Omega: Z^k\to\mathcal{W}$, the reward function is
$\mathcal{R}: \mathcal{S}\times\mathcal{A}\to\mathbb{R}$, $\gamma\in[0,1)$ is
the discount factor, and $\mu_0$ is the initial-state distribution.

At each step, the controller receives the pair $(o_t,w_t)$ with
$o_t\sim\mathcal{H}(x_t)$ and
$w_t = \Omega(z_{t-k+1:t})$ computed deterministically from the recent
telemetry history.  A reliability-aware policy is therefore
$\pi:\mathcal{O}\times\mathcal{W}\to\Delta(\mathcal{A})$, while a
quality-ignorant policy is the restricted map
$\bar{\pi}:\mathcal{O}\to\Delta(\mathcal{A})$.
\end{definition}

\begin{remark}
The episode semantics are
$x_0\sim\mu_0$,
$z_t\leftarrow \mathrm{telemetry}(x_t,\mathrm{fault}_t)$,
$w_t\leftarrow\Omega(z_{t-k+1:t})$,
$o_t\leftarrow \mathrm{parse}(z_t)$,
$a_t\sim\pi(o_t,w_t)$,
$x_{t+1}\sim\mathcal{T}(x_t,a_t)$,
and $r_t=\mathcal{R}(x_t,a_t)$.  Safety is measured on the true-state path
$x_{0:T}$, not on the observation path.  Standard MDPs are recovered when
$\mathcal{O}=\mathcal{S}$, $\mathcal{H}(x)=\delta_x$, and $w_t\equiv 1$.
Classical POMDPs are recovered when the observation channel is partial but
observation quality is not itself exposed to the controller as a runtime
state variable.
\end{remark}

\section{The Hidden Safety Gap}
\label{sec:hidden-gap}

We now state the core safety problem in the battery domain.

\begin{definition}[Battery-domain OASG]
\label{def:battery-oasg}
Let $\mathrm{SOC}_t^{\mathrm{true}}$ denote the physical state of charge
and $\widehat{\mathrm{SOC}}_t$ the value available to the controller.
The \emph{battery-domain Observation-Action Safety Gap} at time~$t$ is the
event
\[
  \mathcal{G}_t^{\mathrm{bat}} \;=\;
  \Bigl\{\,a_t \in \mathcal{A}\!\bigl(\widehat{\mathrm{SOC}}_t\bigr)\Bigr\}
  \;\cap\;
  \Bigl\{\,a_t \notin \mathcal{A}\!\bigl(\mathrm{SOC}_t^{\mathrm{true}}\bigr)\Bigr\},
\]
where $\mathcal{A}(s)$ is the set of dispatch actions that keep the
subsequent SOC within the feasible range $[\mathrm{SOC}_{\min},
\mathrm{SOC}_{\max}]$.
\end{definition}

When $\mathcal{G}_t^{\mathrm{bat}}$ occurs, the controller has committed
a \emph{physically unsafe} action while believing it to be safe.

\section{State Objects and Telemetry Model}

\subsection{True State}

The true state at time~$t$ is the vector
\[
  x_t \;=\; \bigl(\mathrm{SOC}_t^{\mathrm{true}},\;
  L_t^{\mathrm{true}},\; S_t^{\mathrm{true}},\; W_t^{\mathrm{true}}\bigr)
  \;\in\; \mathcal{X} \subset \mathbb{R}^4,
\]
comprising the physical SOC and the true load, solar, and wind power at
the grid connection point.

\subsection{Observed State}

The observed state is the output of the telemetry channel:
\[
  \hat{x}_t \;=\; \mathcal{T}(x_t,\, \xi_t),
\]
where $\mathcal{T}$ is a (possibly stochastic) telemetry operator and
$\xi_t$ encodes the fault realisation at time~$t$ (dropout mask, delay
value, spike indicator, staleness flag).

\subsection{Observation Quality Score}
\label{sec:oqe-formula}

The Observation Quality Engine (OQE) computes a scalar reliability
score from the telemetry features:
\[
  w_t \;=\; \max\!\bigl(w_{\min},\;
    p_{\mathrm{miss},t} \cdot p_{\mathrm{delay},t}
    \cdot p_{\mathrm{ooo},t} \cdot p_{\mathrm{spike},t}\bigr),
\]
where each penalty factor $p_{\cdot} \in [0, 1]$ reflects a specific
fault channel:

\begin{description}
  \item[$p_{\mathrm{miss},t}$:] Fraction of expected readings that
    arrived in the current window.  Full reception yields $p_{\mathrm{miss}} = 1$;
    complete dropout yields $p_{\mathrm{miss}} = 0$.
  \item[$p_{\mathrm{delay},t}$:] Penalty for readings whose timestamp
    exceeds the freshness threshold: $p_{\mathrm{delay}} =
    \exp(-\lambda_d \cdot \bar{\tau}_t)$, where $\bar{\tau}_t$ is the
    mean delay in the window and $\lambda_d$ is a configurable decay rate.
  \item[$p_{\mathrm{ooo},t}$:] Fraction of readings arriving in
    chronological order.  Out-of-order arrivals indicate network jitter
    or replay.
  \item[$p_{\mathrm{spike},t}$:] Penalty for readings flagged as
    transient spikes: $p_{\mathrm{spike}} = 1 - n_{\mathrm{spike}} /
    n_{\mathrm{total}}$.
\end{description}

The floor $w_{\min} > 0$ ensures that the uncertainty inflator always
produces a finite interval; the default value is $w_{\min} = 0.05$.

\section{Fault Model}
\label{sec:fault-model}

We define the fault model as a tuple
$\mathcal{F} = (\rho, \bar{\tau}, \pi_s, \pi_{\mathrm{stale}}, \sigma_d)$:

\begin{itemize}
  \item $\rho \in [0, 1]$: dropout rate (probability that a packet is
    missing at any given step).
  \item $\bar{\tau} \ge 0$: mean packet delay in multiples of the
    decision interval $\Delta t$.
  \item $\pi_s \in [0, 1]$: probability that a reading is a transient
    spike.
  \item $\pi_{\mathrm{stale}} \in [0, 1]$: probability that a sensor
    output is frozen (stale).
  \item $\sigma_d > 0$: standard deviation of the covariate drift between
    calibration and deployment distributions.
\end{itemize}

The CPSBench harness (Chapter~\ref{ch:orius-bench-battery}) instantiates
each fault channel independently and provides composite scenarios
(e.g., \texttt{drift\_combo}) that activate all channels simultaneously.

\section{Battery Dynamics and Feasible Action Set}
\label{sec:battery-dynamics}

\subsection{SOC Dynamics}

The true SOC evolves according to:
\begin{equation}
\label{eq:soc-dynamics}
  \mathrm{SOC}_{t+1}^{\mathrm{true}} \;=\;
  \mathrm{SOC}_t^{\mathrm{true}}
  \;+\; \frac{\eta_c\,[a_t]^{+} \;-\; [a_t]^{-}/\eta_d}{E_{\max}}
  \;\cdot\; \Delta t,
\end{equation}
where:
\begin{itemize}
  \item $a_t \in \mathbb{R}$ is the dispatch action (positive = charge,
    negative = discharge),
  \item $[a_t]^{+} = \max(a_t, 0)$ and $[a_t]^{-} = \max(-a_t, 0)$
    denote the charge and discharge components,
  \item $\eta_c, \eta_d \in (0, 1]$ are the charge and discharge
    round-trip efficiencies,
  \item $E_{\max}$ is the battery energy capacity (MWh), and
  \item $\Delta t$ is the decision interval (hours).
\end{itemize}

\subsection{Feasible Action Set}

Given a SOC value~$s$, the feasible action set is:
\begin{equation}
\label{eq:feasible-set}
  \mathcal{A}(s) \;=\; \Bigl\{\,a \in [-\bar{P},\, \bar{P}] \;:\;
  \mathrm{SOC}_{\min} \;\le\;
  s + \frac{\eta_c [a]^{+} - [a]^{-}/\eta_d}{E_{\max}} \cdot \Delta t
  \;\le\; \mathrm{SOC}_{\max}\Bigr\},
\end{equation}
where $\bar{P}$ is the rated power limit (MW) and
$[\mathrm{SOC}_{\min}, \mathrm{SOC}_{\max}]$ are the regulatory SOC
bounds (typically 20\%--80\% for grid-scale lithium-ion systems).

\section{Dispatch Objective}

The upstream optimiser maximises expected arbitrage revenue over a
rolling horizon~$H$:
\begin{equation}
\label{eq:dispatch-obj}
  \max_{a_t, \ldots, a_{t+H-1}} \;\sum_{\tau=t}^{t+H-1}
  \pi_\tau \cdot a_\tau \cdot \Delta t
  \quad \text{s.t.} \quad
  a_\tau \in \mathcal{A}\!\bigl(\widehat{\mathrm{SOC}}_\tau\bigr)
  \;\;\forall\, \tau,
\end{equation}
where $\pi_\tau$ is the electricity price at step~$\tau$.  Note that the
constraints are evaluated on the \emph{observed} SOC; the optimiser has
no access to the true state.  The DC3S safety layer post-processes the
solution to enforce constraints on the true state via uncertainty
inflation.

\section{The Calibration Gap Under Degradation}
\label{sec:calibration-gap}

Conformal prediction intervals are calibrated on a held-out calibration
set collected under nominal telemetry conditions.  When deployment
telemetry degrades, the residual distribution shifts:

\begin{enumerate}
  \item Dropout introduces a positive bias in the imputed SOC (stale
    holds tend to overestimate SOC during discharge and underestimate
    during charge).
  \item Delay increases the variance of the residual because the
    controller acts on outdated information.
  \item Covariate drift invalidates the exchangeability assumption that
    underpins the conformal guarantee.
\end{enumerate}

The RUI stage compensates by scaling the conformal half-width by
$1/w_t$:
\begin{equation}
\label{eq:rui-inflation}
  q_{\mathrm{eff},t} \;=\; \frac{q_{\alpha}}{w_t},
\end{equation}
where $q_{\alpha}$ is the nominal conformal quantile.  As $w_t$
decreases, the effective half-width grows, widening the uncertainty set
and tightening the safe-action filter.

\section{Baseline Policies}
\label{sec:baselines}

We compare DC3S against three baseline dispatch policies:

\begin{description}
  \item[Deterministic LP.]  Solves~\eqref{eq:dispatch-obj} directly on
    $\hat{x}_t$ with no uncertainty handling.
  \item[Robust MPC.]  Maintains a fixed uncertainty tube
    $[\hat{x}_t - \epsilon, \hat{x}_t + \epsilon]$ with hand-tuned
    margin~$\epsilon$.  The tube does not adapt to observation quality.
  \item[Stochastic risk-aware.]  Samples $N$ scenarios from historical
    forecast errors and enforces CVaR constraints.  The scenario
    distribution is not conditioned on telemetry quality.
\end{description}

All baselines operate on the observed state and do not maintain a
true-state trajectory.  Their safety performance is therefore evaluated
\emph{ex post} by the CPSBench dual-trajectory harness.

\section{Research Questions Formalised}
\label{sec:rqs-formal}

We restate the research questions of Chapter~\ref{ch:introduction} in
mathematical terms.

\begin{description}
  \item[\textbf{RQ1} (Existence of the OASG).]
    Under fault model $\mathcal{F}$ with dropout rate $\rho > 0$, does
    there exist a time step~$t$ such that the OASG
    event~$\mathcal{G}_t^{\mathrm{bat}}$ (Definition~\ref{def:battery-oasg})
    occurs for at least one baseline policy?
  \item[\textbf{RQ2} (Closing the gap).]
    Does DC3S achieve
    $\mathrm{TSVR} = \frac{1}{T}\sum_{t=1}^T
    \mathbf{1}[\mathcal{G}_t^{\mathrm{bat}}] = 0$
    for all tested fault configurations, and at what intervention
    rate~$\mathrm{IR}$?
  \item[\textbf{RQ3} (Formal guarantee).]
    Under the bounded battery-domain assumption set A1--A5 collected in
    Appendix~\ref{app:assumptions}, does the coverage
    inequality
    \[
      \Pr\!\bigl[\,\mathrm{SOC}_{t+1}^{\mathrm{true}}
      \in [\mathrm{SOC}_{\min}, \mathrm{SOC}_{\max}]\,\bigr]
      \;\ge\; 1 - \alpha
    \]
    hold for all $t$ under fault model $\mathcal{F}$?
\end{description}

\section{Scope Boundary}
\label{sec:scope}

The formal results in this dissertation are stated for the following
scope:

\begin{itemize}
  \item \textbf{Domain.}  Grid-scale lithium-ion battery energy storage
    (100~MWh / 50~MW reference plant).
  \item \textbf{Decision granularity.}  Hourly dispatch ($\Delta t = 1$~h).
  \item \textbf{Fault model.}  The five-channel model of
    Section~\ref{sec:fault-model}, with dropout rates up to $\rho = 0.40$.
  \item \textbf{Forecast horizon.}  Up to 48~hours (the maximum horizon
    evaluated in the 48-hour trace studies).
  \item \textbf{Plant model.}  Linear SOC dynamics~\eqref{eq:soc-dynamics}
    with constant efficiencies.  Degradation effects (capacity fade,
    resistance growth) are not modelled in the core pipeline but are
    examined in the extension chapters.
\end{itemize}

Results outside this scope (sub-hourly dispatch, non-lithium chemistries,
vehicle-to-grid applications) require re-calibration and are deferred to
future work.

\section{Dataset Scope}
\label{sec:dataset-scope}

All experiments use two public data sources:

\begin{enumerate}
  \item \textbf{EIA-930 (US).}  Hourly demand, net generation, and
    interchange data for US balancing authorities.  We use the CISO
    (California ISO) and MISO (Midcontinent ISO) regions as primary test
    beds, with PJM, NYISO, and ERCOT for transfer-learning studies.
  \item \textbf{ENTSO-E (Germany).}  Hourly load, solar, and wind data
    from the German transmission system operators.  Used for
    cross-geography generalisation experiments.
\end{enumerate}

Synthetic telemetry faults are injected by the CPSBench harness
(Chapter~\ref{ch:orius-bench-battery}) on top of the real data,
preserving the statistical properties of the underlying load and renewable
signals.

\section{Chapter Summary}

This chapter has defined the ORIUS system architecture, the
five-channel telemetry fault model, the SOC dynamics, the dispatch
objective, and the observation quality engine.  The research questions
have been formalised in terms of the OASG event and the true-state
violation rate.  The scope boundary and dataset scope delimit the claims
made in subsequent chapters.  Chapter~\ref{ch:safety-illusion} uses these
definitions to demonstrate the observed-state safety illusion and motivate
the dual-trajectory evaluation paradigm.
